\documentclass[12pt,a4paper]{article}
\usepackage[utf8]{inputenc}
\usepackage[spanish]{babel}
\usepackage{graphicx}
\usepackage{hyperref}
\usepackage{geometry}
\usepackage{xcolor}
\usepackage{listings}
\usepackage{fancyhdr}
\usepackage{titlesec}

% Configuracion de Margenes
\geometry{left=2.5cm, right=2.5cm, top=3cm, bottom=3cm}

% Colores Personalizados
\definecolor{samsungblue}{HTML}{034EA2}
\definecolor{lightgray}{gray}{0.95}

% Configuracion de Codigo
\lstset{
    backgroundcolor=\color{lightgray},
    basicstyle=\ttfamily\small,
    breaklines=true,
    captionpos=b,
    commentstyle=\color{green!60!black},
    keywordstyle=\color{blue},
    numbers=left,
    numberstyle=\tiny\color{gray},
    frame=single,
}

% Encabezado y Pie de Pagina
\pagestyle{fancy}
\fancyhf{}
\rhead{\color{samsungblue}Retail Vision AI}
\lhead{Reporte Técnico - Equipo Bisontitos}
\cfoot{\thepage}

% --- Portada ---
\begin{document}

\begin{titlepage}
    \centering
    \vspace*{1.5cm}
    {\Huge\bfseries\color{samsungblue} Retail Vision AI \par}
    \vspace{0.5cm}
    {\Large\itshape Sistema de Inventario Inteligente mediante Vision Computacional \par}
    \vspace{2cm}
    {\large Presentado por: \par}
    \vspace{0.3cm}
    {\Huge\bfseries Equipo Bisontitos \par}
    \vspace{2cm}
    \vfill
    {\large Samsung Innovation Campus \par}
    {\large 2026 \par}
\end{titlepage}

\tableofcontents
\newpage

\section{Introducción}
El proyecto \textbf{Retail Vision AI} es una solución integral diseñada para la automatización del control de inventarios en el sector retail. Utilizando tecnologías de vanguardia como \textit{Edge Computing} y \textit{Vision Artificial}, el sistema permite el monitoreo constante de estantes, la detección precisa de productos y la generación de alertas en tiempo real.

El núcleo del sistema reside en una \textbf{Raspberry Pi 5} equipada con un \textbf{AI HAT Hailo-8L}, lo que permite realizar inferencias de redes neuronales profundas con un bajo consumo energético y alta velocidad (hasta 13 TOPS).

\section{Arquitectura del Sistema}
El flujo de datos del proyecto se divide en cuatro capas principales:
\begin{enumerate}
    \item \textbf{Capa de Captura:} Cámaras PTZ y Domo distribuidas estratégicamente en la tienda.
    \item \textbf{Capa de Procesamiento (Edge):} Raspberry Pi 5 que gestiona la captura y el preprocesamiento de imágenes.
    \item \textbf{Capa de Inteligencia Artificial:} Modelos YOLOv8/v11 para detección, ResNet para clasificación y algoritmos de Machine Learning para análisis.
    \item \textbf{Capa de Visualización:} Interfaz web dinámica construida con Three.js que muestra el estado de la tienda en 3D.
\end{enumerate}

\section{Dataset SKU-110K}
Para el entrenamiento y validación de los modelos de detección en ambientes de alta densidad (donde los productos están muy pegados), se utilizó el dataset \textbf{SKU-110K}. 
\begin{itemize}
    \item \textbf{Densidad:} Promedio de 147 objetos por imagen.
    \item \textbf{Escala:} Más de 11,000 imágenes con 1.7 millones de bounding boxes.
    \item \textbf{Desafío:} El modelo debe ser capaz de distinguir entre productos idénticos en estantes saturados, minimizando los falsos negativos.
\end{itemize}

\section{Implementación de Modelos}

\subsection{Detección de Objetos (YOLOv8/World)}
Se utiliza \textbf{YOLOv8x} (y versiones experimentales de YOLO-World v2) para localizar cada producto en el estante. El sistema implementa un \textit{Ensemble} con \textbf{WBF (Weighted Boxes Fusion)} para combinar las predicciones de distintos modelos, logrando precisiones cercanas al 98.5\% mAP en condiciones controladas.

\subsection{Clasificación de Productos (SKU Identification)}
Una vez detectado el objeto, el sistema utiliza un modelo de \textbf{Clasificación} basado en \textbf{ResNet18/50}. Mediante \textit{Transfer Learning}, se entrena al modelo para identificar SKUs específicos (ej: Cereales, Bebidas, Snacks) a partir de los recortes (\textit{crops}) generados por el detector.

\subsection{Regresión de Stock}
Debido a que las cámaras solo ven la "primera fila" de productos, se implementó un modelo de \textbf{Regresión} (Random Forest o XGBoost). Este modelo toma variables como:
\begin{itemize}
    \item Conteo de bounding boxes.
    \item Área total ocupada.
    \item Profundidad estimada del estante.
\end{itemize}
Con esto, el sistema estima cuántas unidades hay realmente en el estante, incluyendo las que están detrás de la primera fila.

\subsection{Clustering de Comportamiento}
Se aplica \textbf{K-Means Clustering} para agrupar los SKUs según su métrica de rotación:
\begin{itemize}
    \item \textbf{Cluster 0:} Alta rotación (requiere reabastecimiento diario).
    \item \textbf{Cluster 1:} Estable (flujo de venta normal).
    \item \textbf{Cluster 2:} Baja rotación (stock inmovilizado).
\end{itemize}

\section{Interfaz de Usuario (3D Demo)}
La interfaz web utiliza un estilo \textbf{Dark Glassmorphism} para ofrecer una experiencia premium. El motor \textbf{Three.js} renderiza una tienda virtual donde el usuario puede:
\begin{itemize}
    \item Ver las cámaras en tiempo real.
    \item Hacer clic en estantes para ver las detecciones AI.
    \item Monitorear el estado de la Raspberry Pi (CPU, Temp, NPU).
    \item Consultar el inventario automatizado.
\end{itemize}

\section{Conclusiones}
El proyecto \textbf{Retail Vision AI} del \textbf{Equipo Bisontitos} demuestra que es posible implementar sistemas industriales de alta complejidad utilizando hardware accesible y optimizado. La combinación de modelos de detección densa con técnicas tradicionales de Machine Learning garantiza un sistema robusto, escalable y eficiente para el futuro del comercio inteligente.

\end{document}
